\chapter{Resumen ejecutivo}
\label{cap:resumen-ejecutivo}

\section*{Prop\'osito de BIOPET}
BIOPET es un sistema web de gesti\'on veterinaria desarrollado como Proyecto
Fin de Curso de la asignatura \asignatura. Permite el registro de usuarios,
el inicio de sesi\'on, y la administraci\'on (alta, consulta, actualizaci\'on
y baja l\'ogica) de mascotas asociadas a un due\~no, con control de acceso
diferenciado por rol (\texttt{ROLE_ADMIN}, \texttt{ROLE_VETERINARIO},
\texttt{ROLE_AUXILIAR}, \texttt{ROLE_DUENO}). El sistema no gestiona datos
cl\'inicos reales ni historiales m\'edicos de mascotas: su alcance actual
cubre identificaci\'on, autenticaci\'on y un cat\'alogo b\'asico de mascotas
por due\~no.

\section*{Arquitectura general}
El sistema sigue una arquitectura de tres niveles, documentada formalmente
mediante el modelo C4 (cap\'itulo~\ref{cap:arquitectura-c4}): un
\emph{frontend} Angular~17 (aplicaci\'on de p\'agina \'unica, componentes
\emph{standalone}), un \emph{backend} Spring~Boot~3.2 sobre Java~21 que
expone una API REST, una base de datos relacional PostgreSQL~16 gestionada
con Flyway, y Redis~7 usado tanto para la lista negra de tokens JWT
revocados como para la cach\'e de lecturas de mascotas. Todo el sistema se
levanta mediante Docker Compose y un \texttt{Makefile} con un \'unico
comando reproducible (\texttt{make up}), sin pasos manuales adicionales.

\section*{Principales mejoras de la Tercera Entrega}
Respecto a la Entrega~1B, la Tercera Entrega incorpor\'o, con evidencia
verificable en el repositorio:
\begin{itemize}[nosep]
  \item Migraci\'on completa del flujo de autenticaci\'on de tokens en el
    cuerpo de la respuesta a cookies \texttt{HttpOnly}+\texttt{Secure}+
    \texttt{SameSite=Strict}, con revocaci\'on real v\'ia Redis, rate
    limiting de login, auditor\'ia estructurada y cabeceras de seguridad
    (Jaime; cap\'itulo~\ref{cap:resultados-jaime}).
  \item HTTPS/TLS~1.3 nativo en el backend, con certificado autofirmado
    acad\'emico y verificaci\'on en vivo del protocolo y el cifrado
    negociados (Jaime).
  \item Separaci\'on de privilegios en PostgreSQL (cuenta de aplicaci\'on
    \texttt{biopet_app} sin permisos de DDL), pinning de im\'agenes Docker
    por digest \texttt{sha256}, y medici\'on de rendimiento y cach\'e con
    k6 y Redis (Fred; cap\'itulo~\ref{cap:resultados-fred}).
  \item Integraci\'on del frontend con el nuevo esquema de cookies, CRUD de
    mascotas accesible (atributos ARIA, manejo de errores anunciado a
    lectores de pantalla) y consolidaci\'on de la ingenier\'ia de
    requisitos (SRS, historias de usuario, casos de uso) (Zaida;
    cap\'itulo~\ref{cap:resultados-zaida}).
  \item Cobertura de pruebas verificada autom\'aticamente con JaCoCo
    (umbral m\'inimo del 60\,\% sobre l\'ineas, ramas y complejidad,
    aplicado a todo el m\'odulo del backend).
\end{itemize}

\section*{Seguridad}
El backend fue auditado contra seis categor\'ias aplicables de OWASP~Top~10
(A01, A02, A03, A05, A07, A09), con evidencia reproducible en
\texttt{docs/mediciones/sec/}: autorizaci\'on por rol y por propietario,
HTTPS/TLS~1.3 real, protecci\'on estructural contra inyecci\'on SQL,
cabeceras de seguridad HTTP, autenticaci\'on por cookies con rate limiting y
revocaci\'on, y auditor\'ia de eventos sin datos sensibles. Las seis
categor\'ias auditadas obtuvieron resultado \textsc{pass}, respaldado por
pruebas automatizadas reales (no solo inspecci\'on manual). El detalle
completo, con limitaciones expl\'icitas por categor\'ia, se presenta en el
cap\'itulo~\ref{cap:resultados-jaime}.

\section*{Reproducibilidad}
El sistema completo se reconstruye desde una clonaci\'on limpia con
\texttt{make up}, sin intervenci\'on manual. Las im\'agenes de terceros
(PostgreSQL, Redis, Maven, Eclipse Temurin) est\'an fijadas por digest
\texttt{sha256} para evitar derivas silenciosas entre reconstrucciones. La
inicializaci\'on de la base de datos usa un mecanismo doble documentado en
ADR-004: Flyway como \'unica fuente de verdad del esquema de aplicaci\'on, y
scripts de \texttt{docker-entrypoint-initdb.d} para crear roles con
privilegios m\'inimos y sembrar datos deterministas.

\section*{Rendimiento}
Se ejecutaron tres corridas de k6 en fr\'io y tres en caliente contra el
endpoint cacheado \texttt{GET /api/mascotas}, sobre HTTPS/TLS~1.3 real. La
tasa de error fue del 0{,}0\,\% en las seis corridas, con un throughput
sostenido de aproximadamente 91--92 peticiones por segundo y una mediana de
latencia de 5{,}7 a 9{,}4~ms seg\'un la corrida. El detalle estad\'istico
completo (intervalos de confianza al 95\,\% con la distribuci\'on t de
Student, percentiles p50/p90/p95/p99) se presenta en el
cap\'itulo~\ref{cap:resultados-fred}.

\section*{Accesibilidad y usabilidad}
El frontend incorpora atributos ARIA (\texttt{aria-invalid},
\texttt{aria-describedby}, \texttt{role="alert"}, \texttt{aria-live}) en los
formularios de autenticaci\'on y en el CRUD de mascotas, verificables
directamente en el c\'odigo fuente. \textbf{Las mediciones cuantitativas de
usabilidad (System Usability Scale) y de accesibilidad automatizada
(Lighthouse) todav\'ia no se han ejecutado a la fecha de este informe}: no
existe en el repositorio ni \texttt{docs/mediciones/sus/} ni
\texttt{docs/mediciones/lighthouse/}. Este informe documenta la
metodolog\'ia prevista para ambas mediciones (cap\'itulo~\ref{cap:protocolo-experimental})
sin reportar resultados que a\'un no existen.

\section*{Resultados generales}
La suite de pruebas del backend ejecuta 109 pruebas reales, con 0 fallos, 0
errores y 0 omitidas (\texttt{BUILD SUCCESS}), verificado en esta misma
fase directamente contra los reportes de Surefire y el reporte JaCoCo
local. La cobertura medida fue del 95{,}87\,\% en l\'ineas, 76{,}87\,\% en
ramas y 80{,}00\,\% en complejidad ciclom\'atica, todas por encima del
umbral autom\'atico configurado del 60\,\%.

\section*{Limitaciones principales}
El certificado TLS es autofirmado y exclusivamente acad\'emico; el rate
limiting de login es en memoria y por instancia, no distribuido; los logs
de auditor\'ia no est\'an integrados con un SIEM; la revocaci\'on de tokens
depende de la disponibilidad de Redis; y toda la evaluaci\'on se realiz\'o
sobre una \'unica instancia del backend en un entorno local. Las mediciones
de usabilidad y accesibilidad automatizada quedan pendientes de ejecuci\'on.
Ninguna de estas limitaciones se presenta como resuelta en este informe; se
detallan con su alcance exacto en el cap\'itulo~\ref{cap:amenazas-validez}.
